%% Extended C++ highlighting, inserted into template_header.tex by `sh gen.sh --new-colors`:
%% keywords blue, types hltype, functions hlfn, numbers hlnum
\definecolor{hltype}{rgb}{0.000,0.502,0.502}
\definecolor{hlfn}{rgb}{0.420,0.208,0.039}
\definecolor{hlnum}{rgb}{0.690,0.000,0.000}
\makeatletter
% function: identifier followed by '(' (peek next char before flushing the identifier)
\let\lst@lparen=(
\let\lst@nextother\relax
\def\lst@AppendOther{%
  \lst@ifletter\expandafter\lst@AppendOtherId\else\expandafter\lst@AppendOtherNoId\fi}
\def\lst@AppendOtherId{\futurelet\lst@nextother\lst@AppendOtherIdA}
% inside a number: '.' and exponent sign after e/E continue it (2.5, 1e-9, 1e+9)
\let\lst@lastletter\relax \let\lst@dot=. \let\lst@plus=+
\expandafter\let\expandafter\lst@che\string e
\expandafter\let\expandafter\lst@chE\string E
\def\lst@ProcessLetter{\lst@whitespacefalse\futurelet\lst@lastletter\lst@AppendLetter}
\def\lst@expsign{\ifx\lst@nextother\lst@umminus\gdef\lst@numcont{1}\fi
  \ifx\lst@nextother\lst@plus\gdef\lst@numcont{1}\fi}
\def\lst@AppendOtherIdA{\gdef\lst@numcont{0}%
  \ifnum\lst@isnum=1
    \ifx\lst@nextother\lst@dot\gdef\lst@numcont{1}\fi
    \ifx\lst@lastletter\lst@che\lst@expsign\fi
    \ifx\lst@lastletter\lst@chE\lst@expsign\fi
  \fi
  \ifnum\lst@numcont=1 \let\lst@nextother\relax\expandafter\lst@Append
  \else\expandafter\lst@AppendOtherIdB\fi}
\def\lst@AppendOtherIdB{\lst@Output\lst@letterfalse\let\lst@nextother\relax
  \lst@AppendOtherNoId}
\def\lst@AppendOtherNoId{\futurelet\lst@lastother\lst@AppendOtherB}
% unary minus: '-' with no operand (identifier, number, ')', ']') before it;
% remember the token without it, in case a digit follows (x=-1, f(-1), return -1)
\expandafter\let\expandafter\lst@umminus\csname lst@um-\endcsname
\let\lst@rparen=) \let\lst@rbrack=]
\gdef\lst@operand{0}\gdef\lst@unary{0}
\newtoks\lst@mtoken \newcount\lst@mlength
\def\lst@AppendOtherB{%
  \ifx\lst@lastother\lst@umminus
    \xdef\lst@unary{\ifnum\lst@operand=0 1\else 0\fi}\gdef\lst@operand{0}%
    \lst@mtoken=\lst@token \lst@mlength=\lst@length
  \else\ifx\lst@lastother\lst@outputspace
  \else\ifx\lst@lastother\lst@rparen \gdef\lst@operand{1}%
  \else\ifx\lst@lastother\lst@rbrack \gdef\lst@operand{1}%
  \else \gdef\lst@operand{0}%
  \fi\fi\fi\fi
  \lst@Append}
% number: a digit not inside an identifier starts a token (with suffixes: 1ll, 1e9, 0x3f)
\gdef\lst@isnum{0}\gdef\lst@curisnum{0}
\def\lst@ProcessDigit{\lst@whitespacefalse\futurelet\lst@lastletter\lst@ProcessDigitA}
\def\lst@ProcessDigitA{%
  \lst@ifletter\expandafter\lst@AppendLetter\else\expandafter\lst@AppendNum\fi}
\def\lst@AppendNum{\gdef\lst@isnum{1}%
  \ifnum\lst@mode=\lst@Pmode\relax\ifx\lst@lastother\lst@umminus\ifnum\lst@unary=1
    \lst@token=\lst@mtoken \lst@length=\lst@mlength \lst@OutputOther
    \lst@lettertrue \expandafter\lst@Append\csname lst@um-\endcsname
  \fi\fi\fi
  \lst@AppendLetter}
\lst@AddToHook{Output}{\xdef\lst@curisnum{\lst@isnum}\gdef\lst@isnum{0}%
  \gdef\lst@operand{0}}
% non-keyword identifier or number; it is an operand for the next '-'
\newcommand\lstidstyle{\gdef\lst@operand{1}%
  \ifnum\lst@curisnum=1 \color{hlnum}%
  \else\ifx\lst@nextother\lst@lparen\color{hlfn}\fi\fi}
\makeatother

\lstdefinelanguage[TRD]{C++}[ISO]{C++}{
  deletekeywords={int,long,short,char,bool,void,auto,unsigned,signed,float,%
    double,dbl,ll,ull,size_t,__int128,string,vector,pair,array,set,map,multiset,%
    unordered_map,unordered_set,queue,deque,priority_queue,stack,bitset,complex,%
    tuple,function,mt19937,mt19937_64,uint32_t,uint64_t,int64_t,uint8_t,uint16_t},
  morekeywords=[1]{nullptr,constexpr,typename,template,using,namespace,operator,%
    inline,static,const,struct,class,true,false,this,new,delete,sizeof,typedef,%
    mutable,public,private},
  morekeywords=[2]{int,long,short,char,bool,void,auto,unsigned,signed,float,%
    double,dbl,ll,ull,size_t,__int128,string,vector,pair,array,set,map,multiset,%
    unordered_map,unordered_set,queue,deque,priority_queue,stack,bitset,complex,%
    tuple,function,mt19937,mt19937_64,uint32_t,uint64_t,int64_t,uint8_t,uint16_t},
  moredirectives={pragma},
}
\lstset{
  defaultdialect=[TRD]C++,
  identifierstyle=\lstidstyle,
  keywordstyle=[1]{\color{blue}},
  keywordstyle=[2]{\color{hltype}},
}
